\section{Open questions (candidate follow-on probes)}

The following five questions are genuinely open after this replication ---
each proposes a concrete numerical or analytical probe that would extend or
stress-test the paper's central graphical-model / tensor-network claim. They
are ordered from most-concrete (noisy extension) to most-speculative
(hierarchical circuit topologies).

\begin{enumerate}
\item \textbf{Noisy-circuit extension.} Does the exact undirected-graphical-model
mapping extend to noisy shallow circuits (depolarizing, amplitude/phase
damping, or realistic hardware noise models), and does the treewidth stay
bounded by $\min(O(d\cdot\ell), O(n))$ once Kraus operators are folded into
the TN? \emph{Next step:} extend \texttt{src/tn\_sim.py} to a density-matrix
TN, verify against full-Liouville-space evolution on $n{\le}10$ at
$p{\in}\{0.001, 0.01, 0.05\}$, and refit the depth slope (expect
$\sim\!2\ell_{\min}$ in Liouville space).

\item \textbf{Head-to-head vs modern classical simulators.} How does the paper's
QuickBB + TensorFlow pipeline compare on identical circuits + hardware
against \texttt{cotengra} hyper-optimized contraction, \texttt{quimb} slicing,
\texttt{qFlex} Feynman-path, and hybrid Schr\"odinger--Feynman? \emph{Next
step:} port $5{\times}9$ depth 12, $7{\times}7$ depth 10, $6{\times}6$ depth
10 circuits to \texttt{quimb+cotengra} with a 10-min hyper-optimizer budget
and report a head-to-head table of largest-intermediate, FLOPs, and
wall-clock.

\item \textbf{ML-boosted approximate inference.} Can complex-valued loopy BP,
neural TN contractors, or GNN-learned elimination-order search beat exact
variable elimination on this regime while keeping controlled error?
\emph{Next step:} implement complex loopy BP on our TN graph and measure
$|\text{BP}-\text{exact}|/|\text{exact}|$ vs iteration on the 70-config
sweep; declare viable only if there is a regime with approx error
$<10^{-3}$ and sub-exponential-in-treewidth cost.

\item \textbf{Sampling-based supremacy-verification primitive.} Can the
graphical-model contraction machinery be reused (shared elimination order,
amortized precomputation of upstream tensors) across many bitstrings for
efficient XEB computation? \emph{Next step:} on our sweep, factor the
network so all output-wire tensors are symbolic and per-bitstring $\langle
x|$ is substituted at the leaves; measure amortized wall-clock per
additional bitstring after the first; extrapolate to $5{\times}9$ depth 12
XEB over $10^5$ samples.

\item \textbf{Non-2D-nearest-neighbor topologies.} Does the
$\min(O(d\ell), O(n))$ bound extend to MERA-like layered ans\"atze,
brickwork+SWAP circuits, or measurement-based cluster-state circuits?
\emph{Next step:} implement MERA on $n{=}16$ with 4 renormalization levels;
implement brickwork+SWAP at densities $\{0.05, 0.15, 0.30\}$; implement 2D
cluster-state MBQC; report the tightest analogous bound that fits each
topology and whether the graphical-model approach is fundamentally tight
there.
\end{enumerate}
